\documentclass[a4paper]{article}
% Español (México) con XeLaTeX: fontspec para fuentes Unicode, polyglossia
% para la división silábica y los nombres del documento en la variante mexicana.
\usepackage{fontspec}
\usepackage{polyglossia}
\setdefaultlanguage[variant=mexican]{spanish}
% Los nombres chinos van como «pinyin (original)»: xeCJK da a los caracteres
% Han su propia fuente; sin ella XeLaTeX los omite con un aviso.
% FandolSong no cubre algunos caracteres de nombres (U+73FA); WenQuanYi Zen Hei sí.
\usepackage[AutoFallBack=true]{xeCJK}
\setCJKmainfont{FandolSong-Regular.otf}
\setCJKfallbackfamilyfont{\CJKrmdefault}{WenQuanYi Zen Hei}
% polyglossia no traduce los nombres de listings.
\AtBeginDocument{\providecommand{\lstlistingname}{}\renewcommand{\lstlistingname}{Listado}%
  \providecommand{\lstlistlistingname}{}\renewcommand{\lstlistlistingname}{Índice de listados}}
\usepackage{amsmath,amssymb}
\usepackage{graphicx}
\usepackage[margin=2.35cm]{geometry}
\usepackage[most]{tcolorbox}
\usepackage{etoolbox}
\usepackage{listings}
\usepackage{xcolor}
\usepackage{hyperref}
\usepackage{booktabs,longtable,tabularx,array}
\usepackage{subcaption}
\usepackage{float}
\usepackage{enumitem}
\usepackage{microtype}
\usepackage{fancyhdr}
\usepackage{needspace}

\definecolor{kbBack}{HTML}{F2F6FA}
\definecolor{kbFrame}{HTML}{9DB4CC}
\definecolor{kbTitle}{HTML}{52708F}
\definecolor{ibBack}{HTML}{FBF5E7}
\definecolor{ibFrame}{HTML}{C9A961}
\definecolor{ibTitle}{HTML}{7D5F1E}
\definecolor{wbBack}{HTML}{FCF2F0}
\definecolor{wbFrame}{HTML}{D6A096}
\definecolor{wbTitle}{HTML}{9E4F43}
\definecolor{dbBack}{HTML}{F4F7F3}
\definecolor{dbFrame}{HTML}{AEC2AC}
\definecolor{dbTitle}{HTML}{5F7F64}

\tcbset{softbox/.style={
  enhanced,breakable,boxrule=0.8pt,arc=2pt,left=3mm,right=3mm,
  top=2mm,bottom=2mm,coltitle=white,fonttitle=\bfseries,
  toptitle=0.6mm,bottomtitle=0.6mm
}}
\newtcolorbox{knowledgebox}[1]{softbox,colback=kbBack,colframe=kbFrame,colbacktitle=kbTitle,title={#1}}
\newtcolorbox{importantbox}[1]{softbox,colback=ibBack,colframe=ibFrame,colbacktitle=ibTitle,title={#1}}
\newtcolorbox{warningbox}[1]{softbox,colback=wbBack,colframe=wbFrame,colbacktitle=wbTitle,title={#1}}
\newtcolorbox{dialoguebox}[1]{softbox,colback=dbBack,colframe=dbFrame,colbacktitle=dbTitle,title={#1}}

\lstset{
  language=Python,basicstyle=\ttfamily\small,
  keywordstyle=\color{kbTitle}\bfseries,stringstyle=\color{wbTitle},
  commentstyle=\color{dbTitle},breaklines=true,frame=single,numbers=left,
  numberstyle=\tiny\color{gray},captionpos=b,columns=fullflexible,
  keepspaces=true,showstringspaces=false,extendedchars=true
}
\BeforeBeginEnvironment{lstlisting}{\Needspace{12\baselineskip}}

\hypersetup{colorlinks=true,linkcolor=kbTitle,urlcolor=kbTitle,citecolor=wbTitle}
\setlist{itemsep=0.25em,topsep=0.4em}
\setlength{\parindent}{2em}
\setlength{\parskip}{0.25em}
\newcolumntype{L}[1]{>{\raggedright\arraybackslash}p{#1}}

\providecommand{\notetitle}{Notas del curso: [tema del curso]}
\providecommand{\noteauthors}{Elaboradas a partir de los materiales públicos de Stanford CS25}
\providecommand{\notedate}{Versión reescrita en 2026}
\providecommand{\videochannel}{Stanford Online}
\providecommand{\videopublishdate}{2022}
\providecommand{\videoduration}{[duración del video]}
\providecommand{\videourl}{}
\providecommand{\videocoverpath}{cover.jpg}
\providecommand{\coursesite}{https://web.stanford.edu/class/cs25/}
\providecommand{\repourl}{https://github.com/hqhq1025/ai-course-notes}
\providecommand{\skillversion}{wdkns/wdkns-skills@39f1a04}

\newcommand{\figsource}[1]{\par\vspace{-0.3em}{\footnotesize\color{gray}\emph{Fuente: #1}}\par}
\newcommand{\teachervoice}[1]{\begin{knowledgebox}{Nota de clase}#1\end{knowledgebox}}
\newcommand{\term}[1]{\textbf{#1}}

\pagestyle{fancy}
\fancyhf{}
\fancyfoot[C]{\thepage}
\fancyfoot[R]{\footnotesize\color{gray}\href{\repourl}{AI Course Notes}}
\renewcommand{\headrulewidth}{0pt}
\renewcommand{\footrulewidth}{0pt}

\newcommand{\makecscover}{
\begin{titlepage}
\centering
\vspace*{0.7cm}
{\Large Stanford CS25 · Transformers United\par}
\vspace{0.8cm}
{\huge\bfseries \notetitle\par}
\vspace{0.6cm}
{\large \noteauthors\par}
\vspace{0.25cm}
{\large \notedate\par}
\vspace{0.9cm}
\ifdefempty{\videocoverpath}{}{
  \includegraphics[width=0.86\textwidth,height=0.43\textheight,keepaspectratio]{\videocoverpath}\par
}
\vfill
\begin{tcolorbox}[width=0.92\textwidth,colback=black!2!white,colframe=black!45,boxrule=0.7pt,arc=2pt]
\textbf{Autor o canal del video}: \videochannel\par
\textbf{Fecha de publicación}: \videopublishdate\par
\textbf{Duración del video}: \videoduration\par
\textbf{Sitio del curso}: \href{\coursesite}{\nolinkurl{\coursesite}}\par
\ifdefempty{\videourl}{}{\textbf{Enlace del video}: \href{\videourl}{\nolinkurl{\videourl}}\par}
\textbf{Normas de elaboración}: \skillversion\ + las normas source-first / teacher-voice / visual-QA de este repositorio
\end{tcolorbox}
\end{titlepage}
\tableofcontents
\newpage
}

\usepackage{multicol}

\renewcommand{\notetitle}{Lecture 36: la llegada de la AGI y los sistemas de Agent utilizables}
\renewcommand{\noteauthors}{Elaborado a partir de la conferencia oficial de Div Garg (AGI, Inc.), la grabación en 1080p y los subtítulos hechos por personas}
\renewcommand{\notedate}{CS25 V5 · versión reescrita source-first de 2026}
\renewcommand{\videochannel}{Stanford Online}
\renewcommand{\videopublishdate}{Clase: 2025-04-15; publicación: 2025-05-13}
\renewcommand{\videoduration}{1:01:01}
\renewcommand{\videourl}{https://www.youtube.com/watch?v=nEHNwdrbfGA}
\renewcommand{\videocoverpath}{cover.jpg}

\newcommand{\lecturefigure}[3]{%
\begin{figure}[H]
  \centering
  \includegraphics[width=0.97\textwidth,height=0.49\textheight,keepaspectratio]{slides-images/#1}
  \caption{#2}
  \figsource{Stanford CS25 V5 official recording, #3}
\end{figure}
}

\let\standardtableofcontents\tableofcontents
\renewcommand{\tableofcontents}{%
  \small
  \setlength{\parskip}{0pt}
  \standardtableofcontents
  \normalsize
  \setlength{\parskip}{0.25em}
}

\begin{document}
\makecscover

\section{Eje del curso: la AGI no es una forma, sino un sistema que puede asumir responsabilidades}

«The Advent of AGI» se presta a leerse como una predicción sobre calendarios, límites de capacidad o superinteligencia. El eje de la clase de Div Garg es más de ingeniería: aunque los modelos ya conversan, razonan y llaman herramientas, \textbf{convertir la inteligencia en un sistema de acción utilizable a diario} todavía exige definir tareas, permisos, memoria, entrenamiento, evaluación, comunicación y recuperación ante fallas. El form factor de la AGI aún no está definido, así que no puede suponerse de antemano que será una ventana de chat más grande; primero hay que preguntar qué responsabilidad le delega el usuario al sistema, en qué se basa la confianza en el sistema y quién puede tomar el control cuando falla.

\lecturefigure{slide-01-title-advent-agi.jpg}{Título del curso: The Advent of AGI}{00:00:15}

\lecturefigure{slide-02-what-does-agi-look-like.jpg}{Pregunta inicial: ¿con qué forma entra la AGI en la vida cotidiana?}{00:01:50}

\teachervoice{00:01:18--00:02:10, el ponente señala que las capacidades de conversación y razonamiento han mejorado rápidamente, pero que todavía no hay una respuesta compartida a «¿cómo se ve la AGI?». Nota de clase: el nivel de capacidad y la forma de producto son dos preguntas distintas; la segunda determina cómo entra el sistema en la sociedad, cómo se verifica y quién asume el riesgo.}

\begin{importantbox}{Definición operativa de esta clase}
Esta clase no define la AGI por una escala de parámetros ni por un umbral de benchmark, sino que la entiende como un \term{límite de responsabilidad} que se amplía gradualmente: el sistema percibe el entorno, formula planes, usa herramientas, mantiene memoria de forma continua y actúa en nombre del usuario. Cada paso que amplía la capacidad debe ampliar al mismo tiempo los mecanismos de evaluación, permisos, observación y recuperación.
\end{importantbox}

\subsection{Los cuatro niveles del contrato de un agent: Memory, Tools, Planning, Action}

El primer diagrama del sistema proviene de la descomposición clásica de Lilian Weng de los LLM-powered agents. Al leer la figura no basta con memorizar los cuatro sustantivos; hay que ver cómo fluye la información: el contexto de corto plazo y la memoria de largo plazo aportan el estado; las herramientas conectan el lenguaje con el mundo exterior; la planeación produce subobjetivos ejecutables; las acciones modifican el entorno y traen de vuelta nuevas observaciones. Si falta cualquiera de estos niveles, el sistema se reduce a un generador de un solo paso incapaz de cerrar el ciclo de forma sostenida.

\lecturefigure{slide-03-agent-architecture-components.jpg}{Arquitectura de un agent: memoria, herramientas, planeación y acción forman en conjunto un ciclo cerrado}{00:03:53}

Un agent basado en herramientas puede formularse como un proceso de decisión parcialmente observable (Partially Observable Markov Decision Process, POMDP):
\[
\mathcal{M}=(\mathcal{S},\mathcal{A},P,R,\Omega,O,\gamma).
\]
Aquí $\mathcal{S}$ es el estado real del entorno, $\mathcal{A}$ es el conjunto de acciones, $P$ es la transición de estados, $R$ es la recompensa, $\Omega$ son las observaciones visibles, $O$ describe cómo el estado genera las observaciones y $\gamma$ es el factor de descuento. Por lo general, el modelo solo ve $o_t\in\Omega$, compuesta por páginas web, mensajes, respuestas de herramientas y memoria comprimida, y no el estado completo $s_t\in\mathcal{S}$; por eso, un siguiente paso que «parece razonable» no es necesariamente la acción correcta a nivel global.

\begin{knowledgebox}{Primer uso: qué resuelve cada uno de Tool, Memory y Planning}
\textbf{Tool} es una capacidad externa que el modelo puede invocar, por ejemplo un navegador, una base de datos, un calendario o un ejecutor de código; \textbf{memory} es el mecanismo para guardar y recuperar el estado entre pasos, lo cual no equivale a meter todo el historial en el prompt; \textbf{planning} consiste en descomponer el objetivo en subobjetivos, revisar el avance y cambiar de ruta tras una falla. Las interfaces entre los tres deben ser explícitas; de lo contrario, no es posible ubicar un error en «el modelo no sabe», «la herramienta falló» o «se perdió el estado».
\end{knowledgebox}

\begin{warningbox}{El riesgo de tomar el CoT como un plan}
Chain-of-Thought (CoT) es una trayectoria de razonamiento expresada en lenguaje y no equivale automáticamente a un plan ejecutable. Un plan real requiere además precondiciones de las acciones, permisos, presupuesto, condiciones de terminación, idempotencia y reversión. Una secuencia de pasos coherente en lo lingüístico puede violar las restricciones del entorno desde el tercer paso.
\end{warningbox}

\subsection{Everyday AGI: la visión debe traducirse de inmediato en una agenda de investigación}

El ponente usa Everyday AGI y la demostración del examen de manejo del DMV para expresar una visión de producto: la IA no solo responde preguntas, sino que puede completar procesos reales en nombre del usuario. La imagen en sí solo demuestra que en ese momento existía una demo de clase; no demuestra que pueda completarse de forma estable en todos los estados, todos los exámenes o todos los sitios web. Más importante aún, la siguiente slide descompone de inmediato la visión en tres líneas de investigación, evaluation, training y communication, lo que indica que la promesa de producto debe convertirse en problemas de sistema verificables.

\lecturefigure{slide-04-everyday-agi.jpg}{Everyday AGI: visión de producto que integra la inteligencia en los procesos cotidianos}{00:04:11}

\lecturefigure{slide-05-dmv-agent-demo.jpg}{Escenario de agente para el examen en línea del DMV mostrado en clase}{00:05:28}

\lecturefigure{slide-06-agent-initiatives.jpg}{De la visión a la agenda de investigación: evaluación, entrenamiento y comunicación entre agents}{00:05:44}

En conjunto, estas tres figuras muestran la progresión «visión--capacidad--infraestructura». Everyday AGI plantea un objetivo perceptible para el usuario, la imagen del DMV ofrece una muestra concreta de interacción y las initiatives enumeran el trabajo de investigación necesario para que esa muestra se convierta en un sistema reproducible. Si solo se conserva el primer nivel, el curso se reduce a una narrativa sobre el futuro; si solo se conserva el tercero, se pierde el porqué de estas técnicas. Un equipo de producto debería escribir cada demo de alto impacto como una especificación de tarea y, a partir de ella, deducir qué deterministic environments, trayectorias de entrenamiento y protocolos entre agents se necesitan.

\teachervoice{00:04:11--00:07:10, el ponente primero presenta la visión de la «AGI cotidiana» y después enumera REAL Evals, AgentQ y agent communication. Las notas advierten: la demo es un punto de entrada para plantear preguntas, no un certificado de confiabilidad; el eje real es cómo evaluar, cómo aprender y cómo colaborar.}

\begin{importantbox}{Cinco preguntas para pasar de la demo a la aceptación del sistema}
\begin{enumerate}
  \item ¿La tarea tiene una condición de finalización que pueda determinarse?
  \item ¿Se puede reproducir la misma prueba después de que cambie el entorno?
  \item ¿Todas las acciones intermedias pueden rastrearse?
  \item ¿Las fallas pueden recuperarse en lugar de reintentarse a ciegas desde el principio?
  \item ¿Las acciones de alto riesgo requieren la confirmación del usuario o la intervención humana?
\end{enumerate}
\end{importantbox}

\subsection{Resumen de la sección}

Esta sección condensa «la llegada de la AGI» en un contrato de sistema: el modelo debe cerrar el ciclo en un entorno parcialmente observable mediante memoria, herramientas, planeación y acción; y toda visión de producto debe traducirse de inmediato en evaluación, entrenamiento, comunicación y límites de responsabilidad. La siguiente sección responde: por qué adoptar una interacción con la computadora «human-like» y cómo se clasifican los niveles de autonomía.
\section{Human-like Agents: la expansión de capacidades debe ir a la par de la jerarquización de permisos}

Esta parte pasa de «por qué se necesitan agents» a «por qué hacer que los agents usen interfaces humanas». La respuesta del ponente no es que la interfaz sea más atractiva, sino que el mundo real ya construyó para las personas navegadores, escritorios, formularios y flujos de inicio de sesión y de pago; si un agent puede operar sobre estas interfaces, puede superar APIs incompletas, servicios heterogéneos y tareas de cola larga. Sin embargo, el control directo también amplía el espacio de acciones y el riesgo; por ello, el nivel de autonomía, los puntos de confirmación y la reversión deben quedar escritos en el contrato del producto.

\lecturefigure{slide-07-why-agents.jpg}{Por qué se necesitan Agents: una sola llamada al modelo no basta para resolver tareas reales}{00:07:10}

\lecturefigure{slide-08-lecture-roadmap.jpg}{Ruta del curso: arquitectura, interacción, memoria, comunicación y direcciones futuras}{00:07:37}

\subsection{De la lista de preguntas a la thesis de interacción}

Las cuatro preguntas «Why, How, Ingredients, What can they do» proporcionan el marco didáctico para lo que sigue. El build final muestra la thesis del ponente: las personas expresan su intención en lenguaje natural y la IA opera la máquina. Lo decisivo aquí no es el lenguaje natural en sí, sino convertir una intención ambigua en acciones acotadas y permitir que el usuario pueda observarlas, modificarlas y deshacerlas.

\lecturefigure{slide-09-building-agents-questions.jpg}{Las cuatro preguntas básicas para construir un Agent}{00:07:54}

\lecturefigure{slide-10-building-agents-thesis.jpg}{Thesis de interacción: la persona usa lenguaje natural y el Agent opera la máquina}{00:08:27}

\lecturefigure{slide-11-ai-agents-why-how.jpg}{Diferencia entre una sola llamada a un foundation-model y un sistema de Agent completo}{00:08:53}

Supongamos que una tarea consta de $n$ pasos que deben tener éxito, y que la probabilidad de éxito del paso $i$ es $p_i$. Si los errores no son recuperables, la tasa de éxito de extremo a extremo es aproximadamente
\[
P(\text{success})=\prod_{i=1}^{n}p_i.
\]
Si cada paso tiene $p_i=0.98$, una tarea de 20 pasos solo alcanza $0.98^{20}\approx 0.668$. Esto explica por qué un modelo que en el chat «se equivoca de vez en cuando» se vuelve muy frágil en cuanto entra en tareas largas; el sistema de agent debe incorporar verificación, reintentos, búsqueda por ramas y recuperación, en lugar de limitarse a esperar que el modelo base mejore un poco más.

\begin{warningbox}{La confiabilidad en tareas largas no es la exactitud de un solo paso}
La exactitud promedio por paso oculta las diferencias entre pasos críticos y sus riesgos. El costo de un error al confirmar un pago, borrar archivos o enviar un mensaje público es mucho mayor que el de una búsqueda inútil; la evaluación de aceptación debe estratificarse según la importancia del paso, la reversibilidad del error y la frecuencia de exposición.
\end{warningbox}

\subsection{Por qué buscar una human-like interface}

La slide completa final presenta cinco motivaciones: reutilizar interfaces diseñadas para personas, convertirse en una extensión digital del usuario, superar las limitaciones de las APIs, usar acciones simples de click/type y aprender de forma continua a partir de la interacción con el usuario. Al leer la figura hay que ver también el reverso: que una interfaz esté diseñada para personas no significa que sea adecuada para una máquina; los cambios visuales, las ventanas emergentes, la publicidad, el inicio de sesión, los mecanismos contra la automatización y la retroalimentación ambigua generan desplazamientos de distribución.

\lecturefigure{slide-12-human-like-agent-benefits.jpg}{Las cinco motivaciones de los Human-like agents y sus riesgos implícitos}{00:11:29}

\teachervoice{00:08:53--00:11:36, el ponente enfatiza que «una sola LLM call no basta» y define al human-like agent como aquel que puede actuar en nombre del usuario sobre las interfaces existentes, gestionar el inicio de sesión y los pagos, usar primitives de click/type y, después, aprender del usuario. Nota de clase: cada ventaja amplía al mismo tiempo los permisos y la superficie de ataque.}

\begin{knowledgebox}{Términos clave: pila de capacidades del Agent}
\begin{tabularx}{\textwidth}{>{\bfseries}p{0.18\textwidth} X X}
Término & Problema que resuelve & Relación con el sistema en esta clase \\
\hline
Browser control & Cómo operar un servicio cuando no hay una API estable & Amplía la cobertura, pero introduce cambios visuales y acciones irreversibles \\
Tool calling & Cómo mapear una solicitud en lenguaje a una llamada estructurada & Es más controlable, pero está sujeto al schema, a los permisos y a los errores devueltos \\
Computer use & Cómo operar el escritorio, el teclado y el mouse, y varias aplicaciones & Las acciones son generales; la observación y la recuperación son más difíciles \\
Long-term memory & Cómo conservar el estado del usuario y de la tarea entre sesiones & Permite continuidad, pero trae riesgos de obsolescencia, privacidad y contaminación \\
Self-correction & Cómo encontrar una ruta alternativa a partir de un estado de falla & Requiere retroalimentación decidible, presupuesto y condiciones de paro \\
\end{tabularx}
\end{knowledgebox}

\subsection{La API y el control directo son dos modelos de riesgo distintos}

La vía de la API ofrece un schema estructurado, códigos de error explícitos y un espacio de acciones más reducido, por lo que suele ser más fácil de auditar; el control directo del navegador o del escritorio cubre las tareas de cola larga, pero la observación son píxeles o el accessibility tree, y el significado de cada acción depende del contexto. Ambas no son «la forma antigua y la nueva», sino dos contracts que deben elegirse según el riesgo de la tarea.

Al leer esta figura comparativa conviene juzgar a lo largo de tres dimensiones. La primera es la \textbf{visibilidad del estado}: la API suele devolver campos estructurados, mientras que el control directo solo puede inferir a partir de la página. La segunda es la \textbf{reversibilidad de la acción}: una interfaz de lectura es fácil de reintentar, pero enviar un formulario o hacer clic en «comprar» puede producir efectos secundarios reales. La tercera es el \textbf{costo de compatibilidad}: las APIs cambian de versión y las páginas web también se reorganizan. Un buen sistema no apuesta de forma permanente por una sola vía: prioriza las interfaces estructuradas, recurre al navegador solo cuando falta cobertura y eleva el nivel de confirmación de las acciones realizadas por esa vía de respaldo.

\lecturefigure{slide-13-agent-computer-two-routes.jpg}{Las dos vías de la Agent--computer interaction: API y control directo}{00:12:31}

\begin{importantbox}{Qué debe registrar explícitamente el contrato de acción}
Para cada acción $a_t$, registre como mínimo el recurso objetivo, los parámetros, quién la invoca, la base de permisos, la clave de idempotencia, la poscondición esperada y el método de reversión. Si la acción es un clic en una página web, también debe convertirse «en qué nodo accesible se hace clic» en un identificador estable, en lugar de guardar solo las coordenadas.
\end{importantbox}

\subsection{Niveles de autonomía: poder hacer algo y deber hacerlo no son la misma pregunta}

El ponente toma prestado el diagrama de niveles de la conducción autónoma para hablar de los agents. Una mayor autonomía no consiste simplemente en reducir los clics del usuario, sino en trasladar de la persona al sistema la responsabilidad de vigilar el entorno, juzgar las anomalías y ejecutar el fallback. Con cada nivel que se sube, hay que explicar quién vigila, quién confirma, quién se encarga de detener la ejecución y si, cuando el sistema falla, sigue existiendo un estado seguro.

\lecturefigure{slide-14-five-levels-autonomy.jpg}{Analogía entre la automatización de la conducción y los niveles de autonomía de un Agent}{00:13:50}

\begin{warningbox}{Límites de la analogía con la conducción autónoma}
Las tareas web y la conducción en carretera comparten la estructura «percepción--decisión--acción--toma de control», pero difieren en riesgo, latencia y responsabilidad legal. La analogía sirve para definir el fallback; no debe usarse para promocionar directamente un producto de agent como «L4/L5».
\end{warningbox}

\subsection{Resumen de la sección}

El valor de los Human-like agents está en reutilizar las interfaces reales y completar flujos de cola larga en nombre del usuario; sus riesgos provienen de esa misma fuente. La API y el control directo deben elegirse según la controlabilidad, la cobertura y la recuperabilidad, y los niveles de autonomía deben diseñarse en sincronía con las responsabilidades de vigilancia, confirmación y toma de control. La siguiente sección usa REAL Bench para mostrar cómo llevar esta complejidad a una evaluación reproducible.
\section{REAL Bench: no reportar solo la puntuación total, sino evaluar la distribución de tareas}

Las páginas web en línea cambian continuamente; si se prueba un agent directamente en sitios reales, aparecen actualizaciones de páginas, cambios de inventario, estados de inicio de sesión y riesgos de seguridad, de modo que los distintos modelos no se comparan en el mismo entorno. El aporte principal de REAL es construir réplicas de sitios de alta fidelidad pero deterministas, para que las trayectorias de acciones puedan ejecutarse de nuevo, los fallos puedan analizarse a posteriori y los modelos puedan compararse sobre el mismo conjunto de tareas.

\lecturefigure{slide-15-real-bench-home.jpg}{REAL: entornos web deterministas de alta fidelidad orientados a Agents}{00:14:43}

\begin{importantbox}{Definición: un entorno determinista no es lo mismo que un entorno simple}
\term{Deterministic simulation} significa que el mismo estado inicial y la misma secuencia de acciones producen resultados reproducibles; aun así, puede incluir páginas complejas, varios sitios y tareas largas. El determinismo resuelve sobre todo la comparabilidad de los experimentos; no garantiza que el entorno cubra todo el ruido del mundo real.
\end{importantbox}

\subsection{De la puntuación total de un modelo a un vector de capacidades por dominio}

La diapositiva del gráfico de anillo muestra la REAL score global de un modelo; el gráfico de radar y el gráfico de barras horizontales revelan las diferencias entre sitios. Al leer la figura, el primer paso no es preguntar «quién quedó primero», sino cómo se construyó el muestreo de tareas, cuántas preguntas tiene cada sitio, si los fallos se concentran en cierto tipo de interacción y si la puntuación global está dominada por tareas sencillas.

Estas tres figuras forman una cadena de diagnóstico que va de lo general a lo particular. La puntuación total sirve para responder «¿el sistema mejoró en conjunto?»; el gráfico de radar por sitio, para responder «¿dónde ocurrió la mejora?», y el gráfico de barras por tarea, para localizar «qué página, herramienta o acción sigue fallando». Si un modelo rinde muy bien en el sitio de compras pero falla de forma persistente en los sitios de calendario y de vuelos, el promedio puede parecer aceptable, pero el equipo de producto debería revisar primero la interpretación de fechas, el estado de identidad, los formularios largos y el cambio entre sitios, en lugar de seguir ampliando el modelo a ciegas.

\lecturefigure{slide-16-real-bench-model-score.jpg}{Puntuación total de los modelos en REAL y estadísticas de tareas por sitio}{00:15:22}

\lecturefigure{slide-17-real-bench-site-radar.jpg}{Perfil de capacidades de distintos Agents en cada sitio}{00:16:16}

\lecturefigure{slide-18-real-bench-task-breakdown.jpg}{Desempeño de los modelos desglosado por sitio web y por tarea}{00:17:22}

Desde la perspectiva de producto, el desglose por sitio y por tarea también determina el orden de las correcciones. Si varios modelos fallan en el mismo sitio web, el problema puede venir del diseño del entorno, de la interfaz de las herramientas o de la ambigüedad de la tarea; si solo un modelo falla en las tareas de formularios largos en todos los sitios, lo más probable es que ese modelo tenga deficiencias en la conservación del estado o en la planificación de acciones. Por eso, una plataforma de evaluación no solo debe producir una tabla de clasificación para la investigación, sino también agrupamientos de fallos accionables, traces reproducibles y contraejemplos mínimos, para que los ingenieros puedan reproducir el problema en lugar de adivinarlo a partir de un video.

Sea $\mathcal{W}$ el conjunto de sitios web, $\mathcal{T}_w$ el conjunto de tareas del sitio $w$ y $s_{w,t}\in\{0,1\}$ el indicador de éxito de una tarea individual. El promedio macro es
\[
S_{\text{macro}}=\frac{1}{|\mathcal{W}|}\sum_{w\in\mathcal{W}}
\frac{1}{|\mathcal{T}_w|}\sum_{t\in\mathcal{T}_w}s_{w,t}.
\]
$S_{\text{macro}}$ da el mismo peso a cada sitio; el promedio micro, en cambio, da más peso a los sitios con más tareas. El reporte debe indicar qué agregación se usa; de lo contrario, el mismo resultado puede contarse como historias distintas.

Además, se dividen las tareas en slices $g\in\mathcal{G}$ según riesgo, longitud y tipo de interacción:
\[
S_g=\mathbb{E}[s\mid g],\qquad
S_{\min}=\min_{g\in\mathcal{G}}S_g.
\]
$S_{\min}$ es el peor slice. Un producto de alto riesgo no puede optimizar solo la media; si las tareas de pago, permisos o eliminación forman el peor grupo, un 90\% global tampoco demuestra que el sistema pueda desplegarse.

\begin{knowledgebox}{Lectura de la figura: qué responde cada una de las tres páginas de resultados}
La página de puntuación total responde «¿en qué nivel se encuentra este sistema en conjunto?»; el gráfico de radar responde «¿las capacidades están equilibradas?», y el gráfico de barras por tarea responde «¿dónde se concentran los fallos?». Las tres deben usarse en conjunto: la puntuación total para ordenar, la distribución para diagnosticar y el peor slice como umbral para la puesta en producción.
\end{knowledgebox}

\begin{warningbox}{Benchmark saturation no equivale a confiabilidad en el mundo abierto}
Las réplicas deterministas permiten pruebas de regresión estables, pero los sitios web reales también incluyen actualizaciones de interfaz, latencia, permisos, personalización, medidas contra la automatización y consecuencias irreversibles. El benchmark debe servir como plataforma de entrenamiento y depuración, no como una forma de convertir directamente la puntuación del leaderboard en una tasa de incidentes reales.
\end{warningbox}

\teachervoice{00:14:03--00:17:58, el ponente subraya que REAL usa réplicas de sitios modernas, complejas y deterministas, y que compara los modelos sitio por sitio. El docente enfatiza: lo esencial de un agent web no es «saber hacer clic en botones», sino poder completar tareas de manera estable en distintos sitios, tareas y estados de error.}

\subsection{Resumen de la sección}

REAL Bench resuelve la reproducibilidad y la visibilidad de la distribución en la evaluación de agents. Un reporte confiable incluye, como mínimo, la puntuación global, el promedio macro por sitio, los slices de tareas, el peor grupo y las causas de fallo a nivel de trayectoria. La siguiente sección aborda AgentQ: cómo lograr que un agent no solo sea evaluado, sino que también aprenda de las trayectorias fallidas.
\section{AgentQ: de la «captura de pantalla del fallo» a trayectorias de recuperación de las que se puede aprender}

Este tramo es la parte de toda la clase que el ritmo de las demo oculta con más facilidad. El ponente muestra primero fallos comunes y después tareas de reservación, calendario y vuelos. El verdadero punto didáctico no es que «el agent puede reservar en un restaurante», sino cómo el sistema busca rutas alternativas, corrige errores, mantiene varias decisiones a la vez y recupera las trayectorias exitosas y fallidas como datos de entrenamiento.

\subsection{Primero, admitir que el Agent no es perfecto}

El collage de fallos incluye validación de formularios, pagos, horarios no disponibles y páginas de error. Le recuerda al lector que la retroalimentación de un entorno real casi nunca es un reward limpio, sino mensajes de error, resultados vacíos, páginas a medio completar y restricciones ocultas. La motivación de investigación de AgentQ es justamente que el agent siga buscando en esos estados, en lugar de dar por terminada la tarea tras un solo fallo.

Esta subsección continúa la distribución de evaluación del capítulo anterior: REAL nos dice en qué sitios se concentran los fallos, y AgentQ pregunta cómo actuar después de fallar. Al leer la figura de apertura, conviene separar primero los errores en «recuperables» y «no recuperables»: un campo faltante en un formulario suele poder corregirse, un pago rechazado puede requerir la intervención del usuario, y una identidad o un destinatario equivocados pueden obligar a terminar la tarea. Solo si se define antes el límite de la recuperación, el algoritmo de búsqueda no confundirá «seguir intentando» con el comportamiento correcto.

\lecturefigure{slide-19-agentq-failure-collage.jpg}{AgentQ abre mostrando tipos de fallo en páginas web reales}{00:17:59}

\lecturefigure{slide-20-agentq-title.jpg}{AgentQ: un marco de búsqueda y aprendizaje para Agent autónomos}{00:18:08}

\begin{importantbox}{El fallo debe estructurarse}
Un fallo se divide, como mínimo, en observation error, planning error, tool/action error, environment constraint y policy/safety rejection. Registrar solo «la tarea falló» no permite decidir si hay que cambiar el prompt, las herramientas, los datos, la recompensa o los permisos.
\end{importantbox}

\subsection{Tarea del restaurante: búsqueda, retroceso y conclusión}

El giro clave de la demo de reservación es que la hora objetivo no está disponible. El Agent consulta primero OpenTable, luego usa Google para obtener información complementaria y finalmente regresa al sitio original para completar la reservación. Los cuatro estados siguientes no están para mostrar una animación, sino para explicar un proceso de recuperación: la acción falla, se busca entre herramientas, se reubica y se completa con confirmación.

\lecturefigure{slide-21-reservation-prompt.jpg}{Objetivo del usuario: restaurante, hora y número de personas especificados}{00:18:11}

\lecturefigure{slide-22-agent-web-actions.jpg}{El Agent ejecuta acciones de reservación en la página web}{00:18:15}

\lecturefigure{slide-23-google-fallback.jpg}{Tras el fallo en el sitio original, recurre a Google para obtener información de disponibilidad}{00:18:19}

\lecturefigure{slide-24-adaptable-site-formats.jpg}{Regresa entre sitios y se adapta a distintas estructuras de página}{00:18:22}

\lecturefigure{slide-25-reservation-confirmed.jpg}{Tarea completada: la confirmación de la reservación se vuelve una poscondición verificable}{00:18:25}

\begin{knowledgebox}{Lectura de la figura: cómo validar la cadena de recuperación}
Primero, confirma si el agent reconoce la restricción del entorno «las 10 PM no están disponibles»; luego, revisa si la búsqueda entre sitios introdujo confusiones con restaurantes homónimos, fechas o zonas horarias; por último, verifica que la página de confirmación coincida con la intención original. La tarea solo cuenta como exitosa si la página final, la orden estructurada y el objetivo del usuario están alineados.
\end{knowledgebox}

\subsection{Self-correction y decisiones paralelas}

El siguiente grupo de figuras desglosa la capacidad en tres niveles: poder corregirse a sí mismo, poder manejar varias decisiones a la vez y poder reconocer una hora de calendario poco razonable y reprogramarla. Aquí, «self-correction» no debe entenderse como que el modelo diga «me equivoqué», sino como una máquina de estados que detecta que falló una poscondición, genera una acción alternativa y conserva las restricciones de la tarea original.

\lecturefigure{slide-26-self-correcting-behavior.jpg}{Árbol de comportamiento del Agent: tras una rama fallida, vuelve a elegir ruta}{00:18:28}

\lecturefigure{slide-27-multiple-decisions.jpg}{Mantener varias decisiones paralelas dentro de una misma tarea}{00:18:31}

\lecturefigure{slide-28-calendar-prompt.jpg}{Tarea de calendario: programar una reunión de una hora mañana a las 3}{00:18:32}

\lecturefigure{slide-29-calendar-rescheduled.jpg}{Reconoce que las 3 AM no son razonables y reprograma a las 3 PM}{00:18:39}

\begin{warningbox}{La «corrección automática» es lo que con más facilidad excede los permisos}
Si el usuario escribe «3», ¿cambiar de 3 AM a 3 PM es una inferencia razonable o una modificación no autorizada? En una tarea de calendario de bajo riesgo puede resolverse con una confirmación; en una tarea de alto riesgo, la corrección propuesta debe presentarse al usuario, no ejecutarse en silencio.
\end{warningbox}

\subsection{Skills routing y contexto de acciones}

«right skills at the right time» enfatiza el enrutamiento de herramientas. La tarea de vuelos requiere interpretar ciudades, fechas, viaje sencillo, preferencias de precio y de asiento, y después elegir la herramienta de búsqueda. El enrutador no es solo un clasificador: también debe decidir qué restricciones entran en los parámetros de la herramienta, cuáles se reservan para el filtrado posterior y si debe cambiar de herramienta tras un fallo.

Los casos anteriores del restaurante y del calendario muestran sobre todo la recuperación; esta sección pasa a la \textbf{selección de capacidades}. Que un agent tenga diez herramientas no significa que deban exponerse todas al modelo al mismo tiempo; cuantas más herramientas hay, más fácil es que ocurran nombres parecidos, parámetros en conflicto y usos indebidos de permisos. Al leer las figuras, conviene fijarse en la fidelidad con la que las restricciones se conservan a lo largo de la cadena: si el usuario dice «preferably one-way» y «window seat», la etapa de búsqueda puede relajar temporalmente el precio, pero no puede olvidar la preferencia de asiento en la recomendación final, y mucho menos cambiar en silencio la fecha de salida porque alguna herramienta respondió más rápido.

\lecturefigure{slide-30-right-skills-right-time.jpg}{Invocar la skill correcta en el momento correcto es un problema de orquestación del Agent}{00:18:41}

\lecturefigure{slide-31-flight-prompt.jpg}{La solicitud de vuelo incluye fecha, viaje sencillo, precio y preferencia de asiento}{00:18:42}

\lecturefigure{slide-32-flight-window-seat.jpg}{En los resultados de búsqueda sigue filtrando asientos de ventanilla en clase económica}{00:18:48}

\lecturefigure{slide-33-booking-success-improvement.jpg}{Mejora en el booking task success rate presentada en clase}{00:18:54}

Esta serie de demo también expone costos ocultos más allá de la tasa de éxito final. La búsqueda entre sitios aumenta la carga de páginas y los token, la corrección del calendario requiere explicar la ambigüedad horaria y el filtrado de vuelos mantiene varias restricciones blandas a la vez; una solicitud aparentemente sencilla de «hazlo por mí» incluye en realidad seguimiento de estado, cambio de herramientas, conservación de preferencias y verificación de resultados. Si el sistema solo optimiza la etiqueta de éxito, puede aprender a hacer clic de forma más agresiva o a ignorar las restricciones blandas. Los datos de entrenamiento deben incluir, como parte del resultado, la fidelidad completa a la intención del usuario.

\teachervoice{00:17:58--00:20:47, el ponente dice explícitamente que AgentQ no es un sistema perfecto; estas demo sirven para mostrar fallos, retrocesos, autocorrección, decisiones múltiples y skill routing. Las notas advierten: los resultados de la clase corresponden a tareas y configuraciones de modelo específicas, y no deben generalizarse como confiabilidad en cualquier sitio web.}

\begin{importantbox}{Una trace mínima reproducible de Agent}
\begin{lstlisting}[caption={Campos mínimos que debe guardar una trayectoria de Agent}]
trace = {
  "goal": user_goal,
  "initial_state": state_snapshot,
  "steps": [
    {"observation": obs, "thought_summary": plan,
     "action": action, "tool_result": result,
     "postcondition": check, "cost": cost}
  ],
  "outcome": outcome,
  "user_confirmations": confirmations,
  "rollback": rollback_record
}
\end{lstlisting}
\end{importantbox}

\subsection{Resumen de la sección}

Las demo de AgentQ no son publicidad de «reservación automática», sino un conjunto de casos de recuperación ante fallos: el entorno impone restricciones, el agent busca rutas alternativas, la máquina de estados verifica las poscondiciones, el enrutador elige la skill y la trayectoria se guarda como material de aprendizaje. El siguiente capítulo formaliza estos fenómenos como MCTS, process feedback y preference learning.
\section{Búsqueda, Critique y Preference Learning: el mecanismo de entrenamiento de AgentQ}

El artículo de AgentQ combina tres mecanismos: Monte Carlo Tree Search (MCTS) amplía la exploración, la LLM critique ofrece retroalimentación de proceso sobre los nodos intermedios y un objetivo al estilo de Direct Preference Optimization (DPO) aprende de las trayectorias mejores y peores. Los tres resuelven, respectivamente, los problemas de «no hay exploración», «no se sabe dónde está el error» y «la experiencia exitosa no se consolida».

\lecturefigure{slide-34-agentq-methods-paper.jpg}{Los tres componentes de AgentQ: MCTS, self-critique/process supervision y DPO}{00:20:47}

\subsection{MCTS: conservar ramas alternativas para tareas largas}

La búsqueda en árbol conserva varias acciones candidatas en un mismo estado, en lugar de comprometerse de forma voraz con la primera ruta. La fórmula de selección UCT típica es

Este paso retoma la autocorrección vista en la demo: si la policy emite una sola acción cada vez, después de un fallo solo puede volver a muestrear desde el mismo contexto, y es fácil que repita la ruta original; la estructura de árbol, en cambio, guarda de forma explícita «qué ya se probó, qué ramas tienen mayor valor y dónde todavía vale la pena explorar». Al leer una figura de MCTS, primero se identifican la root, las ramas candidatas y los nodos hoja, y después se distingue entre la value estimate y el resultado real de la ejecución. Para un agent web, el árbol de búsqueda debe desplegarse de preferencia en un simulador o en una copia que admita reversión; en el sitio real solo se ejecuta un número reducido de acciones ya filtradas.
\[
a^*=\arg\max_a\left[Q(s,a)+c\sqrt{\frac{\ln N(s)}{N(s,a)+1}}\right].
\]
$Q(s,a)$ es la estimación de valor actual, $N(s)$ es el número de visitas al estado padre, $N(s,a)$ es el número de visitas a la acción y $c$ controla la exploration. El primer término favorece las rutas que ya se sabe que son buenas; el segundo incentiva probar las acciones menos visitadas.

\lecturefigure{slide-35-mcts-search.jpg}{MCTS equilibra exploración y explotación en el árbol de decisiones del Agent}{00:21:40}

\begin{warningbox}{MCTS en entornos web no es un retroceso gratuito}
Muchas acciones son irreversibles: enviar un pedido, mandar un mensaje o realizar un pago no se pueden deshacer como en un tablero de ajedrez. Una implementación segura debe buscar en un entorno simulado o desplegar solo estados que admitan reversión; la exploración en el entorno real debe tener límites de permisos, de presupuesto y de sandbox.
\end{warningbox}

\subsection{Process feedback: recompensar las decisiones intermedias, no solo el desenlace}

Usar únicamente el éxito final $r_{\text{outcome}}$ produce una señal dispersa y, además, puede etiquetar como ejemplo positivo una trayectoria que «tuvo éxito por casualidad, pero con un proceso peligroso». La recompensa de proceso puede escribirse como
\[
R(\tau)=\alpha r_{\text{outcome}}+
\sum_{t=1}^{T}\beta_t r_{\text{process},t}-\lambda C(\tau),
\]
$\tau$ es la trayectoria completa, $r_{\text{process},t}$ evalúa si el paso $t$ es razonable y $C(\tau)$ registra el costo de las herramientas, las infracciones o el riesgo irreversible. $\alpha,\beta_t,\lambda$ determinan el equilibrio entre resultado, proceso y costo.

\lecturefigure{slide-36-self-critique-process.jpg}{La LLM critique ofrece retroalimentación de proceso sobre las acciones candidatas}{00:22:58}

\begin{knowledgebox}{Los tres niveles de la critique}
La \textbf{syntax critique} revisa el formato de la acción; la \textbf{state critique} revisa si la acción es coherente con la página actual y sus restricciones; la \textbf{goal critique} revisa si la acción sigue sirviendo al objetivo original del usuario. Una critique que solo evalúa la fluidez del lenguaje no puede detectar errores de fecha, identidad, permisos ni presupuesto.
\end{knowledgebox}

\subsection{Construir preferencias a partir de ramas exitosas y fallidas}

AgentQ ordena las ramas que produce la búsqueda en preferred y rejected, y después actualiza la policy con un objetivo de preferencias fuera de línea. La forma habitual de DPO es
\[
\mathcal{L}_{\text{DPO}}(\theta)
=-\mathbb{E}\log\sigma\!\left(
\beta\left[
\log\frac{\pi_\theta(y^+\mid x)}{\pi_{\text{ref}}(y^+\mid x)}-
\log\frac{\pi_\theta(y^-\mid x)}{\pi_{\text{ref}}(y^-\mid x)}
\right]\right).
\]
$x$ es el estado junto con el contexto de la tarea, $y^+$ es la acción o trayectoria mejor, $y^-$ es la rama peor, $\pi_{\text{ref}}$ es la política de referencia y $\beta$ controla la intensidad de la preferencia. En el caso de los agents, las preferencias no pueden provenir solo de la captura de pantalla final de la página; deben combinarse con la critique de proceso, los errores de las herramientas y las restricciones de seguridad.

\lecturefigure{slide-37-rlhf-preference-learning.jpg}{Construir señales de aprendizaje de preferencias a partir de trayectorias de búsqueda y critique}{00:23:23}

\teachervoice{00:20:47--00:23:24, el ponente explica AgentQ mediante MCTS, LLM critique y preference learning. El docente enfatiza: tanto las trayectorias exitosas como las fallidas deben entrar en el aprendizaje; sobre todo, hay que lograr que el modelo aprenda a hacer backtrack y recover en entornos complejos.}

\subsection{La trayectoria de OpenTable en siete pasos: el éxito final no oculta los errores intermedios}

Este conjunto de figuras debe leerse como una secuencia de transiciones de estado, no como siete capturas parecidas. La trayectoria primero navega hacia el restaurante objetivo, entra a la página principal, localiza el restaurante y luego selecciona una fecha equivocada; después abre el selector de fechas, corrige la fecha y completa la reservación. El error ocurre en el cuarto paso, la recuperación en el quinto y el sexto, y la confirmación final en el séptimo.

El valor didáctico de esta trayectoria es que incluye a la vez la \textbf{conservación del objetivo} y la \textbf{corrección local}. El sistema no se puso a buscar otro restaurante porque la fecha estuviera mal, ni redujo la solicitud del usuario a «basta con reservar en cualquier restaurante»; conservó el restaurante, el número de personas y el rango de horario, y solo modificó el campo erróneo. Al leer cada figura conviene seguir tres líneas: el estado actual de la página, la siguiente acción que el agent declara y las restricciones del objetivo original que todavía no se cumplen. Solo así se puede juzgar si la recuperación es una corrección real o si se llegó por casualidad a un resultado que podía completarse.

\lecturefigure{slide-38-opentable-step1-navigate.jpg}{Paso 1: navegar al restaurante y al sitio objetivo}{00:24:10}

\lecturefigure{slide-39-opentable-step2-homepage.jpg}{Paso 2: entrar a la página principal de OpenTable e iniciar la búsqueda}{00:24:15}

\lecturefigure{slide-40-opentable-step3-restaurant.jpg}{Paso 3: localizar la página de La Pizza \& La Pasta}{00:24:42}

\lecturefigure{slide-41-opentable-step4-wrong-date.jpg}{Paso 4: se seleccionó una fecha equivocada; la trayectoria presenta un fallo diagnosticable}{00:24:49}

\lecturefigure{slide-42-opentable-step5-date-selector.jpg}{Paso 5: volver a abrir el selector de fechas}{00:24:52}

\lecturefigure{slide-43-opentable-step6-correct-date.jpg}{Paso 6: corregir la fecha y entrar al flujo de confirmación}{00:24:54}

\lecturefigure{slide-44-opentable-step7-confirmed.jpg}{Paso 7: la reservación se completa; la página de confirmación aporta la evidencia final}{00:24:56}

\begin{importantbox}{Trajectory-level acceptance}
Registrar el éxito final como $s_T=1$ todavía no basta. Se recomienda reportar al mismo tiempo: el número de acciones erróneas, el número de recuperaciones, el número de acciones irreversibles, el número de confirmaciones del usuario, el costo total de las herramientas y la latencia hasta completar la tarea. Solo así se distingue entre «completarse sin contratiempos a la primera» y «dar rodeos, hacer clics erróneos, pero acabar teniendo éxito por casualidad».
\end{importantbox}

\begin{lstlisting}[caption={Pseudocódigo de verificación y recuperación de trayectorias}]
for step in range(max_steps):
    observation = environment.observe()
    action = policy.propose(goal, observation, memory)
    if not permission_check(action):
        action = request_user_confirmation(action)
    result = environment.execute(action, idempotency_key=step)
    if postcondition_holds(goal, result):
        return SUCCESS
    if repeated_state(result) or budget_exhausted():
        return HUMAN_FALLBACK
    memory.record(observation, action, result)
return TIMEOUT
\end{lstlisting}

\subsection{Página de resultados: las cifras deben ir ligadas a la configuración experimental}

La gráfica de tasa de éxito en OpenTable que se mostró en clase compara varias configuraciones y destaca la mejora que aportan componentes como AgentQ y online search. El artículo reporta que Llama-3 70B sin ejemplos pasa de 18.6\% a 81.7\%, y que al agregar online search alcanza 95.4\%. Estas cifras son resultados bajo una tarea, una recolección de datos, un modelo y un método de evaluación específicos; respaldan que «la búsqueda y el aprendizaje son eficaces», pero no que «todas las tareas web ya son confiables al 95\%».

Al leer la gráfica de barras, primero se compara el cambio al agregar componentes sobre el mismo base model y después las diferencias entre modelos distintos; no hay que atribuir a un solo algoritmo la diferencia entre todas las barras. También hay que revisar si el criterio de éxito considera solo la página final, si se permite búsqueda humana o externa, cuántos datos de interacción usó cada configuración y si los fallos provienen de la misma distribución de tareas. Si las configuraciones experimentales difieren, la altura de las barras solo puede tomarse como el resultado de la combinación del sistema; de ella no se puede deducir directamente la contribución causal independiente de un solo módulo de MCTS, DPO o search.

\lecturefigure{slide-45-agentq-real-world-benchmark.jpg}{Comparación de la tasa de éxito de AgentQ en escenarios reales de reservación}{00:26:41}

\begin{warningbox}{No hay que confundir la mejora relativa con la confiabilidad absoluta}
Pasar de 18.6\% a 81.7\% es un enorme avance de investigación, pero 81.7\% significa que aproximadamente una de cada cinco tareas sigue fallando. Incluso con 95.4\%, si se ejecutan cien acciones de alto riesgo al día, el número esperado de fallos sigue siendo inaceptable. El umbral para llevarlo a producción debe considerar la frecuencia, la gravedad y la posibilidad de recuperación.
\end{warningbox}

\subsection{Resumen de la sección}

La cadena de mecanismos de AgentQ es la siguiente: la búsqueda en árbol genera ramas diversas, la process critique evalúa los pasos intermedios, el objetivo de preferencias reincorpora a la política las ramas buenas y malas, y la evaluación a nivel de trayectoria revisa a la vez el resultado y el proceso. El caso de OpenTable muestra que los errores no son lo preocupante; los defectos sistémicos son no poder detectarlos, no poder recuperarse de ellos y no poder aprender de ellos.
\section{Neural Compute, Memory y Personalization}

A continuación, el ponente compara al agent con un sistema de cómputo: el modelo funciona como una neural compute unit invocable, las invocaciones repetidas forman un ciclo, el scratchpad guarda el estado de trabajo y la memoria a largo plazo asume la persistencia entre sesiones. La analogía tiene mucho valor didáctico, pero hay que dejar claro que el Transformer no es una CPU ni el embedding store es un disco. La analogía sirve para separar el cómputo, la memoria de trabajo y el estado persistente.

\subsection{El modelo como Neural Compute Unit}

La primera figura asigna la secuencia máxima de tokens de entrada a la secuencia máxima de tokens de salida y subraya que una invocación del modelo es un cómputo dentro de una ventana finita. La segunda, la analogía con las instrucciones MIPS, indica lo siguiente: un procesador tradicional tiene una semántica de instrucciones fija, mientras que la «instrucción» de un modelo de lenguaje se interpreta a partir del lenguaje natural y del contexto en conjunto; por eso es flexible, pero inestable.

\lecturefigure{slide-46-neural-compute-unit.jpg}{Una invocación del modelo abstraída como neural compute unit}{00:28:00}

\lecturefigure{slide-47-mips-instruction-analogy.jpg}{Analogía entre las instrucciones fijas de MIPS y la invocación de un modelo de lenguaje natural}{00:28:34}

\begin{knowledgebox}{Correspondencias de la analogía}
Los tokens de entrada equivalen al programa y los datos; los parámetros del modelo, a una estructura de cómputo fija; los tokens de salida, al resultado del cómputo, y el scratchpad, al área de trabajo. Sin embargo, el modelo de lenguaje carece de la semántica estricta, las transacciones y el determinismo de una ISA de hardware; por ello, cada invocación necesita schema, validación y manejo de errores que lo compensen.
\end{knowledgebox}

\subsection{Looped Transformer: convertir una generación en un cómputo iterativo}

La figura del Looped Transformer devuelve la salida a la entrada y forma un ciclo entre scratchpad, memory e instructions. Los sistemas de agents hacen precisamente esto: mediante orchestration, extienden una sola generación hacia adelante a un algoritmo de varios pasos. Si el estado interno en la iteración $k$ es $z_k$, puede abstraerse como
\[
z_{k+1}=F_\theta(z_k,o_k,m_k),\qquad
a_k=G(z_{k+1}),
\]
donde $F_\theta$ es el cómputo del modelo, $o_k$ es la nueva observación, $m_k$ es la memoria recuperada y $G$ asigna el estado interno a una acción de herramienta. El sistema debe definir cuándo detenerse; de lo contrario, el propio ciclo se convierte en un modo de falla.

\lecturefigure{slide-48-looped-transformer.jpg}{Looped Transformer: scratchpad, memory e instructions entran en el circuito iterativo}{00:28:51}

\begin{importantbox}{La condición de paro importa más que «pensarlo un poco más»}
Las condiciones de terminación incluyen, como mínimo, la tarea completada, el presupuesto agotado, el estado repetido, el umbral de riesgo, la cancelación por parte del usuario y la intervención humana. Un reasoning loop sin condición de paro convierte la incertidumbre en tokens, latencia y costos de herramientas.
\end{importantbox}

\subsection{Long-term Memory: persistencia, recuperación, actualización y olvido}

La slide de memoria a largo plazo describe la memory como long-lived y persistent, y enumera embedding, retrieval models, jerarquía, consistencia temporal, estructura y adaptación en línea. La primera vez que aparecen hay que distinguir: la \term{working memory} es el estado de la tarea actual; la \term{long-term memory} es la información persistente entre sesiones; la \term{model memory} puede referirse al conocimiento aprendido en los parámetros. Las tres tienen mecanismos de actualización distintos.

\lecturefigure{slide-49-long-term-memory.jpg}{Mecanismos y problemas abiertos de la memoria a largo plazo}{00:35:17}

Una expresión básica de recuperación es
\[
m_t=\operatorname{TopK}_{i}\;\operatorname{sim}(q_t,e_i),
\]
donde $q_t$ es el embedding de la consulta actual, $e_i$ es la representación de la entrada de memoria y $m_t$ son las top-$k$ entradas devueltas. Recuperar solo por similitud trae contenido caduco o sensible; por eso, la puntuación real debe incorporar tiempo, fuente, permisos y confianza:
\[
\text{score}_i=\alpha\,\text{sim}(q_t,e_i)+\beta\,\text{recency}_i+
\eta\,\text{trust}_i-\lambda\,\text{privacyRisk}_i.
\]

\begin{warningbox}{Contaminación de la memoria y persistencia de errores}
Si un malentendido de una conversación se escribe en la memoria a largo plazo, se reforzará una y otra vez en el futuro. Antes de escribir, hay que distinguir entre lo que el usuario declaró explícitamente, lo que el sistema infirió y el estado temporal de la tarea; las preferencias de alto impacto deben poder ser consultadas, corregidas y eliminadas por el usuario.
\end{warningbox}

\begin{lstlisting}[caption={Registro de memoria con fuente y política de caducidad}]
memory_item = {
  "fact": normalized_content,
  "source": source_event,
  "confidence": confidence,
  "scope": ["travel", "restaurants"],
  "created_at": timestamp,
  "expires_at": expiry,
  "sensitivity": sensitivity,
  "user_editable": True
}
\end{lstlisting}

\subsection{La Personalization es un problema de alineación usuario--Agent}

La slide de personalización enumera a la vez preferencias explícitas e implícitas. Las explícitas, como alergias, asientos y platillos favoritos, suelen poder confirmarse directamente; las implícitas provienen de estadísticas de comportamiento y es más fácil que confundan una elección fortuita, una restricción de precio o un contexto histórico con un gusto estable. El objetivo del sistema puede escribirse como una optimización con restricciones:
\[
\max_\pi\;\mathbb{E}[U_{\text{user}}(\tau)]
\quad\text{s.t.}\quad
C_{\text{privacy}}(\tau)\le \epsilon_p,
\;C_{\text{safety}}(\tau)\le \epsilon_s.
\]
$U_{\text{user}}$ es la utilidad para el usuario, $C_{\text{privacy}}$ y $C_{\text{safety}}$ son los costos de privacidad y seguridad, y $\epsilon_p,\epsilon_s$ son los umbrales permitidos. La personalización no puede lograrse a costa de recopilar datos más allá de los límites.

\lecturefigure{slide-50-personalization.jpg}{Personalización: las preferencias explícitas e implícitas constituyen juntas la user-agent alignment}{00:36:42}

\lecturefigure{slide-51-personalization-challenges.jpg}{Retos de la personalización: recopilación de datos, aprendizaje, adaptación en línea y privacidad}{00:37:56}

Las dos figuras de personalización deben leerse junto con la de memoria a largo plazo: la memory decide «qué guardar», la personalization decide «cómo usarlo» y la privacy decide «cuándo no puede usarse». Por ejemplo, si un usuario eligió un asiento de pasillo para un viaje de negocios, eso no debe generalizarse de forma permanente como preferencia global; el sistema debe guardar la fuente, el contexto y la confianza, y volver a confirmar antes de una decisión de alto impacto. La personalización auténtica no consiste en que el agent adivine más, sino en que sepa qué preferencias son confiables, cuáles caducaron y cuáles deben preguntarse de nuevo.

\teachervoice{00:35:17--00:38:00, el ponente compara la memoria a largo plazo con un disco, define la personalización como user-agent alignment y enfatiza la consulta activa, el aprendizaje pasivo, SFT, human feedback, on-the-fly adaptation y privacy. Nota de clase: las preferencias aprendidas deben llevar fuente y un mecanismo para revocarlas.}

\begin{knowledgebox}{Términos clave: Memory y Personalization}
\begin{tabularx}{\textwidth}{>{\bfseries}p{0.2\textwidth} X X}
Término & Mecanismo central & Falla principal \\
\hline
Embedding memory & Recuperación de textos o eventos por similitud & Semánticamente cercano pero factualmente erróneo; caducidad temporal \\
Episodic memory & Guarda interacciones concretas y sus resultados & Mucho ruido; sensible en privacidad \\
Semantic memory & Condensa muchos eventos en hechos estables & Los errores de generalización quedan fijados a largo plazo \\
Online adaptation & Cambia la conducta al instante según la nueva retroalimentación & Deriva catastrófica; manipulación mediante retroalimentación maliciosa \\
Preference model & Predice las elecciones o la satisfacción del usuario & Confunde elecciones históricas con valores reales \\
\end{tabularx}
\end{knowledgebox}

\subsection{Resumen de la sección}

La analogía neural-compute ayuda a separar en capas la invocación del modelo, el ciclo, el estado de trabajo y la memoria a largo plazo, pero no sustituye el diseño real de interfaces. La memoria a largo plazo requiere recuperación, fuente, caducidad, permisos y edición por parte del usuario; la personalización requiere aprender bajo restricciones de utilidad, privacidad y seguridad, no recopilar datos de comportamiento sin límite.
\section{Multi-Agent Systems: paralelización, especialización y protocolos de comunicación}

Ante tareas largas, un solo agent está limitado por el contexto, el número de herramientas y la acumulación de errores. Los sistemas multi-agent intentan dividir la tarea entre distintos worker y dejar que un manager coordine los resultados. Así se puede paralelizar y especializar, pero también aparecen nuevos problemas de sistemas distribuidos: la división de tareas, la consistencia de los mensajes, la ejecución duplicada, la espera de sincronización, la propagación de permisos y la recuperación ante fallos.

\subsection{Por qué usar varios Agents}

La primera figura representa los sistemas autónomos con varios robots. El build completo enumera parallelization y task specialization. La paralelización solo reduce la latencia cuando los subtareas dependen poco entre sí. Si todos los worker esperan el mismo estado compartido, el costo de comunicación anula la ganancia.

Esta sección pasa de la memory y la personalization de un solo agent a los límites de la colaboración. Antes de dividir una tarea conviene trazar su grafo de dependencias. Consultar varias fuentes independientes puede hacerse en paralelo, pero redactar el informe final tiene que esperar a que se reúna la evidencia. Cada worker puede usar sus propias herramientas especializadas, pero la identidad del usuario y los permisos de pago, que son compartidos, requieren control central. Al leer la figura no hay que tomar el número de robots como medida de capacidad. Hay que preguntar si cada agent tiene una responsabilidad independiente, si sus salidas pueden combinarse, quién resuelve los conflictos y si el manager aísla a un worker que excede el tiempo límite o devuelve contenido malicioso.

\lecturefigure{slide-52-multi-agent-autonomous-systems.jpg}{Forma básica de un sistema autónomo multi-Agent}{00:39:49}

\lecturefigure{slide-53-why-multiagent-systems.jpg}{Las dos motivaciones principales de los sistemas multi-Agent: paralelización y especialización}{00:40:48}

Si la tarea se divide en $K$ subtareas, la latencia en serie es $\sum_k T_k$. La latencia en paralelo ideal es aproximadamente
\[
T_{\text{parallel}}\approx \max_k T_k+T_{\text{coord}}+T_{\text{merge}}.
\]
$T_{\text{coord}}$ es el tiempo de planificación y comunicación, y $T_{\text{merge}}$ es el de combinación y verificación de resultados. Si el costo de coordinación es demasiado alto o si la subtarea más larga domina, un sistema multi-agent no es más rápido por sí solo.

\begin{warningbox}{«Multi-Agent» no significa copiar el mismo Prompt varias veces}
Un sistema multi-agent sin límites de rol, sin protocolo de estado compartido y sin resolución de conflictos por lo general solo aumenta el costo en tokens. Cada worker debe tener una entrada definida, un schema de salida, permisos, un plazo límite y un manejo de fallos.
\end{warningbox}

\subsection{Hierarchy y syncing primitives}

El diagrama jerárquico entrega la solicitud del usuario a un manager, que luego la reparte entre varios worker. La responsabilidad del manager no se limita a reenviar mensajes: también planifica las dependencias, controla el presupuesto, decide cuando los resultados entran en conflicto y reasigna el trabajo si un worker falla. Las llamadas \term{syncing primitives} son los mecanismos que permiten que varios agent lleguen a un acuerdo operativo sobre el estado compartido, por ejemplo locks, números de versión, barrier, confirmaciones de mensajes e ID de tarea idempotentes.

\lecturefigure{slide-54-agent-hierarchy.jpg}{Jerarquía manager--worker y flujo de información entre Agents}{00:42:21}

\begin{lstlisting}[caption={Contrato de tarea manager--worker}]
task = {
  "task_id": stable_id,
  "goal": scoped_goal,
  "inputs": immutable_inputs,
  "permissions": allowed_tools,
  "depends_on": prerequisite_task_ids,
  "deadline": deadline,
  "output_schema": schema,
  "retry_policy": retry_policy
}
\end{lstlisting}

\subsection{MCP y A2A: un estándar de conexión no es una garantía de seguridad}

La slide de protocolos asigna MCP (Model Context Protocol) a la conexión entre modelos, herramientas y datos, y A2A (Agent-to-Agent) a la comunicación entre agents. El valor de MCP está en unificar la interfaz de descubrimiento e invocación, de modo que cada aplicación no tenga que escribir glue code a la medida para cada herramienta. A2A, en cambio, necesita definir la identidad, los mensajes, el estado de las tareas y la semántica de la colaboración. Ambos resuelven la interoperabilidad, pero no resuelven por sí solos la autorización ni la corrección.

\lecturefigure{slide-55-mcp-a2a-protocols.jpg}{MCP conecta modelos y herramientas; A2A permite la colaboración entre Agents}{00:43:44}

\begin{knowledgebox}{Primer uso: Protocol, Schema y Authorization}
\textbf{Protocol} define cómo se comunican las partes; \textbf{schema} define los campos de los mensajes y sus tipos; \textbf{authorization} determina quién puede ejecutar qué acción. Que un MCP server exponga una herramienta no significa que cualquier client deba tener permiso para invocarla. Los permisos tienen que ser mínimos, contar con capacidad de auditoría y poder revocarse.
\end{knowledgebox}

\subsection{Por qué se compara MCP con USB-C}

La primera figura de MCP destaca la conexión unificada entre modelos, datos y herramientas. El build final enumera además resilience, dynamic tool discovery e intent-based execution. Al leer la figura, «absorber los cambios de la API» debe entenderse como aislamiento mediante una capa de abstracción, no como una interfaz que nunca falla: los cambios en el schema del server, en la semántica y en los permisos siguen requiriendo gestión de versiones y pruebas de compatibilidad.

\lecturefigure{slide-56-mcp-usbc.jpg}{La analogía de MCP con USB-C: conexión unificada entre modelos y herramientas}{00:45:11}

\lecturefigure{slide-57-mcp-api-benefits.jpg}{Resiliencia y descubrimiento dinámico de MCP frente a la integración punto a punto de API}{00:45:46}

La estandarización del protocolo también cambia la localización de fallos. Cuando falla una integración punto a punto, las responsabilidades de la aplicación, del modelo y de la API suelen mezclarse. Con una capa de protocolo se puede revisar por separado si el server es descubrible, si el schema coincide, si los permisos bastan, si la invocación excedió el tiempo límite y si la respuesta pasó la validación. El costo es que la propia capa de protocolo se vuelve una dependencia crítica: hay que versionarla, monitorearla e impedir que un server malicioso inyecte descripciones de herramientas engañosas. La llamada resilience proviene de una separación clara en capas y de las pruebas, no de un nombre unificado.

\teachervoice{00:38:00--00:46:09. El ponente resume el valor de los sistemas multi-agent en parallelization y specialization, y enfatiza hierarchy, synchronization, MCP, A2A, resilient integrations y dynamic discovery. Nota de clase: el lenguaje natural es ambiguo por naturaleza, así que la comunicación necesita un protocolo estructurado y un estado.}

\begin{importantbox}{Lista mínima de seguridad para la capa de protocolo}
El descubrimiento de herramientas debe separarse de la concesión de permisos. Las herramientas sensibles requieren un scope explícito. Cada invocación incluye la identidad de quien llama y el ID de la tarea. Las salidas requieren validación de schema. Las actualizaciones del protocolo deben pasar pruebas de compatibilidad. Las acciones de alto riesgo deben requerir la confirmación del usuario o un control doble.
\end{importantbox}

\subsection{Resumen de la sección}

Los sistemas multi-agent convierten un problema de un solo modelo en un problema de sistemas distribuidos. La paralelización y la especialización solo aportan beneficios cuando la tarea puede dividirse y la comunicación está bajo control. El contrato manager--worker, las primitives de sincronización, los schema de MCP/A2A, la identidad y los permisos determinan en conjunto si el sistema es confiable.
\section{Deployment Contract: Reliability, Loop, Observabilidad y Human Override}

La última slide didáctica reúne todo lo anterior en cuatro tipos de problemas de despliegue: confiabilidad, ciclos y deriva del plan, pruebas y benchmark, y despliegue real y observabilidad. No es un apéndice, sino el conjunto de condiciones de aceptación de la visión «Everyday AGI». Cuanto más se acerca el sistema al correo electrónico, al calendario, a los pagos y a las cuentas públicas del usuario, menos puede sustituirse la garantía operativa por la tasa de éxito promedio de una demo.

\lecturefigure{slide-58-key-issues-autonomous-agents.jpg}{Problemas clave de los Agent autónomos: confiabilidad, ciclos, pruebas y observabilidad}{00:46:09}

Esta página de síntesis se corresponde con la Everyday AGI de la apertura: aquella pregunta «¿qué pasaría si la inteligencia entrara en la vida cotidiana?», y esta responde «¿qué debe cumplirse antes de que entre?». La confiabilidad determina si se puede otorgar autorización; el loop control determina el costo y el límite de la pérdida de control; el testing determina si un cambio puede verificarse; la observabilidad determina si un incidente puede explicarse. Las cuatro no son parches de operación posteriores al lanzamiento, sino que deben diseñarse desde las etapas de agent architecture, recolección de datos y entorno de entrenamiento.

\subsection{Reliability: incluso con 99.9\% hay que preguntar por la frecuencia de la tarea y el costo del incidente}

Si la probabilidad de fallo de una sola acción de alto riesgo es $q$ y se ejecuta $N$ veces al año, la probabilidad de al menos un fallo es
\[
P(\ge 1\text{ failure})=1-(1-q)^N.
\]
Con $q=0.001$ y $N=1000$, esa probabilidad es de aproximadamente $63.2\%$. Por lo tanto, «99.9\% reliable» no es seguro por sí mismo; además hay que reducir las acciones automáticas de alto riesgo, añadir confirmaciones, limitar los permisos y garantizar que todo sea recuperable.

\teachervoice{00:46:09--00:50:00, el ponente dice que los pagos, la banca, el correo electrónico y las cuentas de redes sociales exigen una confiabilidad cercana al nivel de producción; el sistema no puede enviar contenido equivocado, realizar transacciones erróneas ni caer en ciclos. El docente enfatiza el audit trail, el monitoreo en línea, el human fallback y la toma de control por parte del usuario.}

\begin{warningbox}{La tasa de éxito promedio oculta la cola catastrófica}
Un informe de confiabilidad debe presentar a la vez la tasa de errores graves, la tasa de errores irreversibles, el uso de permisos sin confirmación, el ciclo más largo, el costo máximo por tarea y la tasa de toma de control humana. Reportar solo el éxito promedio de las tareas no demuestra que el sistema no provocará incidentes de baja frecuencia y alto daño.
\end{warningbox}

\subsection{Loop y plan divergence: dar al sistema un presupuesto y un estado detectable}

Un ciclo puede consistir en repetir la misma acción, o en un plan que se expande sin cesar sin avanzar. La detección puede combinar el hash del estado, los n-gram de acciones, la distancia al objetivo y el presupuesto:
\[
\text{loopRisk}_t=
\alpha\,\mathbf{1}[h(s_t)=h(s_{t-k})]
+\beta\,\text{actionRepeat}_t
+\gamma\,\text{noProgress}_t.
\]
Cuando el riesgo supera el umbral, el sistema debe detenerse, degradarse o solicitar intervención humana, en lugar de aumentar automáticamente el token budget.

\begin{lstlisting}[caption={Ciclo de ejecución con presupuesto, detección de repeticiones y toma de control humana}]
while not done:
    if spent_cost > cost_limit or elapsed > time_limit:
        return escalate("budget exceeded")
    if state_hash in recent_states:
        return escalate("loop detected")
    action = agent.next_action()
    if action.risk >= confirmation_threshold:
        action = wait_for_user_confirmation(action)
    execute(action)
    append_audit_log(action)
\end{lstlisting}

\subsection{Q\&A: de zero-shot a task-specific training}

Ante la pregunta «¿cómo pasar del 40\% a casi el 99\%?», el ponente señala que muchos frontier models operan en zero-shot frente a las agent interfaces: la distribución de preentrenamiento no incluye las páginas web ni los flujos de operación específicos. Mediante RL específico de la tarea, autocorrección y recolección repetida de trayectorias, un dominio estrecho puede elevarse; pero eso no equivale a que todas las tareas del mundo abierto se saturen al mismo tiempo.

\teachervoice{00:50:00--00:54:20, el ponente usa AgentQ para mostrar que el task-specific RL puede aumentar notablemente la exactitud en tareas de reservación, y a la vez reconoce que los sitios nuevos, las tareas nuevas y las restricciones nuevas siguen produciendo desplazamiento de distribución. La nota advierte: la saturación local y la confiabilidad general deben reportarse por separado.}

El ciclo virtuoso del despliegue puede escribirse como
\[
\begin{aligned}
\text{production traces}&\rightarrow\text{failure taxonomy}
\rightarrow\text{offline simulation}\\
&\rightarrow\text{training}\rightarrow\text{shadow eval}
\rightarrow\text{controlled rollout}.
\end{aligned}
\]
En cada paso deben conservarse las versiones y la provenance de los datos; solo así puede determinarse si una mejora proviene del modelo, del prompt, de las herramientas o de un cambio en el entorno.

\subsection{Hallucination y domain-specific regression suites}

La respuesta de ingeniería del ponente tiene dos niveles: un modelo base más potente reduce algunos errores; cada producto concreto sigue necesitando sus propias pruebas y evaluaciones. Para un agent vertical, debe mantenerse una regression suite formada por tareas reales, entradas extremas, límites de permisos, fallos de herramientas e incidentes históricos, y reproducirse a diario después de cualquier cambio en el modelo, el prompt, el schema de las herramientas o la policy.

\teachervoice{00:54:20--00:57:55, el ponente enfatiza construir alrededor de mil domain scenarios relevantes, revisar de forma continua los cambios de prompt y las regresiones, comparar distintos modelos y optimizar aún más con fine-tuning o reinforcement learning. Experiencia práctica: no tomar la actualización del modelo del proveedor como la aceptación propia.}

\begin{importantbox}{Un umbral de liberación entregable para un Agent}
\begin{enumerate}
  \item la strict offline suite no presenta regresiones graves;
  \item el shadow traffic no ejecuta efectos secundarios reales;
  \item se habilitan primero los usuarios canary y las herramientas de bajos permisos;
  \item las acciones de alto riesgo deben confirmarse y poder revertirse;
  \item el monitoreo cubre éxito, costo, ciclos, errores y toma de control humana;
  \item un incidente puede rastrearse hasta la versión del modelo, prompt, memory, tool o policy.
\end{enumerate}
\end{importantbox}

\subsection{Modelos pequeños, Manager--Worker y la prueba del mundo real}

La sesión de Q\&A discute smaller versus larger models. El ponente considera que un modelo pequeño entrenado con reasoning traces y RL puede ser más eficaz en tareas estrechas, y propone una posible arquitectura con un modelo grande como manager y modelos más pequeños como workers. Aquí debe respetarse el límite de la evidencia: el tamaño del modelo es solo una de las variables de enrutamiento; también cuentan la calidad, la latencia, el costo, la privacidad, el uso de herramientas y las necesidades de contexto.

\teachervoice{00:57:55--00:59:47, el ponente considera que se necesitan modelos «más inteligentes, no solo más grandes», y propone la posible combinación large manager + small workers; al mismo tiempo, toma el real-world task success como litmus test, en lugar de decidir según el número de parámetros.}

Si una tarea $x$ puede enrutarse a un modelo $m$, un objetivo sencillo es
\[
m^*(x)=\arg\max_m\left[
Q(m,x)-\lambda_L L(m,x)-\lambda_C C(m,x)-\lambda_R R(m,x)
\right],
\]
donde $Q$ es la calidad, $L$ la latencia, $C$ el costo y $R$ el riesgo. El enrutador debe elegir dinámicamente según la tarea, en lugar de fijar de manera permanente que «el manager debe ser el más grande y el worker el más pequeño».

\subsection{El problema de la memory no tiene una única respuesta de hardware}

La última pregunta compara la agent memory con la RAM, la ROM y el hard drive. El ponente no ofrece una arquitectura unificada, sino que enfatiza que cada aplicación es distinta y combina componentes distintos. Esta reserva es muy importante: algunas tareas solo necesitan un scratchpad de corto plazo, otras requieren un perfil de usuario editable y otras necesitan registros de eventos y una base de conocimiento; llamar «memory» a todo esto conduce a compartir información indebidamente y a una persistencia excesiva.

\teachervoice{00:59:47--01:00:57, ante la analogía RAM/ROM/disk, el ponente afirma explícitamente que no hay una straight answer; hay que buscar los modelos y componentes adecuados según la aplicación, como coding, chat o actions. Nota de clase: la memory architecture es una decisión de producto y de riesgo; más no siempre es mejor.}

\subsection{Resumen de la sección}

El contrato de despliegue convierte los indicadores de investigación del Agent en responsabilidades reales: la confiabilidad debe considerar la frecuencia de las tareas, los ciclos requieren detección y presupuesto, las pruebas deben cubrir la distribución del dominio, y en producción debe haber trace, monitoreo, auditoría y toma de control humana. El RL local puede mejorar tareas de dominio estrecho, pero no elimina el desplazamiento de distribución; tanto el tamaño del modelo como la estructura de la memory deben enrutarse según la tarea.
\section{Síntesis y ampliación}

Esta clase parte de la forma abstracta de la AGI y termina en un mapa de sistemas que se puede ejecutar. Lo esencial de un Agent no es que «el modelo sepa invocar herramientas», sino que objetivos, estado, acciones, permisos, evaluación, aprendizaje, memoria, comunicación y recuperación formen un ciclo cerrado. Si falta cualquiera de estas capas, una demostración de capacidades se convierte en una automatización imposible de mantener.

Además, esta clase separa con claridad la «inteligencia del modelo» de la «confiabilidad del sistema»: la primera determina el límite superior de las acciones candidatas; la segunda determina qué acciones pueden llegar al mundo real. Un modelo más potente puede reducir algunos errores de razonamiento, pero no aporta por sí mismo auditoría, autorización, reversión, control de versiones ni responsables. Al estudiar sistemas de Agent, toda mejora del modelo debe evaluarse siempre dentro del plano de control completo.

\subsection{Una tabla que articula todo el Agent stack}

\begin{center}
\begin{tabularx}{\textwidth}{>{\bfseries}p{0.16\textwidth} X X X}
Capa & Pregunta central & Mecanismos representativos & Señales de falla \\
\hline
Goal & Qué quiere realmente el usuario & Análisis de restricciones, confirmación & Deriva del objetivo, pérdida de restricciones implícitas \\
Policy & Qué hacer en el siguiente paso & planning, MCTS, routing & Comportamiento voraz, ciclos, exceso de permisos \\
Environment & Dónde se ejecutan las acciones & browser/API/tools & Cambios en la página, acciones irreversibles \\
Evaluation & Cómo saber que se hizo bien & REAL, slice, trace checks & La puntuación global oculta la cola \\
Learning & Cómo mejorar a partir de las trayectorias & critique, DPO, RL & Explotación de la recompensa, datos contaminados \\
Memory & Cómo conservar el estado entre pasos y sesiones & scratchpad, retrieval, profile & Caducidad, filtraciones, persistencia de errores \\
Coordination & Cómo colaboran varios Agent & hierarchy, sync, MCP/A2A & Conflictos, duplicación, esperas \\
Deployment & Cómo asumir responsabilidad en el mundo real & budget, audit, override & Pérdida de control, falta de trazabilidad, imposibilidad de recuperación \\
\end{tabularx}
\end{center}

\begin{importantbox}{La síntesis más importante de esta clase}
La forma de producto de la AGI aún no está definida, pero las restricciones de ingeniería de un Agent confiable ya son claras: \textbf{evaluar de forma reproducible, aprender sobre trayectorias, actuar según permisos, recordar con fuentes, colaborar mediante protocolos y, ante una falla, detenerse y devolver el control.}
\end{importantbox}

\subsection{Doce preguntas antes de iniciar un proyecto de Agent}

\begin{multicols}{2}
\small
\begin{enumerate}
  \item ¿Puede una máquina determinar el objetivo del usuario y la condición de término?
  \item ¿Qué acciones son reversibles y cuáles requieren confirmación?
  \item ¿El Agent actúa mediante una API o controla directamente la interfaz?
  \item ¿El benchmark cubre sitios web reales, la longitud de las tareas y los slice de riesgo?
  \item ¿Se registra la trace completa de observation--action--result?
  \item ¿La taxonomy de fallas permite ubicar el origen en el modelo, las herramientas, la memory o la policy?
  \item ¿La búsqueda y la self-critique tienen costo y condición de paro?
  \item ¿Los preference data incluyen tanto ramas exitosas como fallidas?
  \item ¿La memoria de largo plazo tiene fuente, caducidad, permisos y edición por parte del usuario?
  \item ¿El sistema multi agent realmente puede paralelizarse o solo agrega comunicación?
  \item Más allá de MCP/A2A, ¿cómo se gestionan la identidad, la autorización y las versiones?
  \item Tras el lanzamiento, ¿quién monitorea, quién toma el control y cómo se revierte?
\end{enumerate}
\normalsize
\end{multicols}

\subsection{Lecturas adicionales}

\begin{itemize}
  \item Putta et al., \emph{Agent Q: Advanced Reasoning and Learning for Autonomous AI Agents}: MCTS, self-critique y aprendizaje por preferencias. \href{https://arxiv.org/abs/2408.07199}{arXiv:2408.07199}
  \item AGI, Inc., \emph{REAL: Benchmarking Autonomous Agents on Deterministic Simulations of Real Websites}: entornos web deterministas y leaderboard. \href{https://github.com/agi-inc/real}{REAL repository}
  \item Lilian Weng, \emph{LLM Powered Autonomous Agents}: panorama de la arquitectura de memory, planning, tools y action. \href{https://lilianweng.github.io/posts/2023-06-23-agent/}{article}
  \item Documentación oficial de Model Context Protocol: protocolo abierto de conexión para herramientas, recursos y prompt. \href{https://modelcontextprotocol.io/}{official docs}
\end{itemize}

\paragraph{Alcance de la evidencia.}
Esta clase registra el contenido impartido en abril de 2025. REAL, AgentQ, MultiOn, MCP y los resultados de los modelos de ese momento deben entenderse según la fecha de la clase; estas notas conservan sus mecanismos y sus lecciones de ingeniería, sin convertir las demo, el estado de los productos ni las cifras de benchmark en garantías vigentes de productos en 2026.

\end{document}
